\documentclass[twocolumn]{article}
\usepackage[utf8]{inputenc}
\usepackage{amsmath}
\usepackage{amssymb}
\usepackage{graphicx}
\usepackage{booktabs}
\usepackage{cite}
\usepackage{hyperref}

\title{Comprehensive Performance Analysis of N-Node Networks under Malicious Attacks and Timeout Constraints}
\author{Antigravity AI Assistant}
\date{May 1, 2026}

\begin{document}

\maketitle

\begin{abstract}
This paper presents an extensive theoretical and empirical investigation into the performance of network architectures under adversarial conditions. We begin with a foundational analysis of single-node models under two distinct attack vectors: destruction and modification. We then extend these models to complex multi-hop environments, including tandem feedforward chains and symmetric feedback networks. By combining renewal-reward theory with moment-matching approximations (Gamma distribution), we provide a robust framework for predicting sojourn time and throughput. Our theoretical models are validated against discrete-event simulations, showing high fidelity in capturing the non-linear delay escalations near system stability limits.
\end{abstract}

\section{Introduction}
Modern network infrastructures are increasingly vulnerable to both localized attacks and systemic congestion. This work reconciles the performance gap between idealized theoretical models and empirical simulation results for N-node systems.

\section{Single-Node Baseline Models}
Before analyzing multi-hop networks, we establish baseline performance for a single node under two attack modalities.

\subsection{Destruction Attacks}
In a destruction attack, the adversary removes packets from the queue at rate $a \lambda_a$. The probability of a packet being lost to an attack before completing service is modeled as:
\begin{equation}
L_{dest} = \frac{a \lambda_a}{\mu + \lambda_a}
\end{equation}
The effective traffic rate $\Lambda^*$ is found by solving:
\begin{equation}
\Lambda^* = \frac{\lambda_n}{1 - (L_{dest} + (1-L_{dest})e^{-T(\mu - \Lambda^*)})}
\end{equation}

\subsection{Modification Attacks}
Modification attacks involve the corruption of packets during service. The loss probability is proportional to the service time:
\begin{equation}
L_{mod} = \frac{a \lambda_a}{\mu}
\end{equation}
The packet consumes service resources at the node even if it is eventually found to be corrupted and requires retransmission.

\section{Multi-Node Network Architectures}

\subsection{Feedforward (Tandem) Network}
In a tandem network of $N$ nodes, each failure requires retransmission from node 0. The journey time for $i$ successful hops follows a hypoexponential distribution. To handle the timeout probability $P(S_i > W)$, we approximate $S_i$ as a Gamma distribution by matching the first two moments:
\begin{equation}
k_{gamma} = \frac{(E[S_i])^2}{Var(S_i)}, \quad \theta_{gamma} = \frac{Var(S_i)}{E[S_i]}
\end{equation}

\subsection{Feedback (Mesh) Network}
In a symmetric feedback network, the traffic rate $\Lambda^*$ is identical for all nodes due to the $1/(N+1)$ transition probability. The effective load is solved via fixed-point iteration:
\begin{equation}
\Lambda^* = \lambda \frac{\sum_{k=1}^{\infty} (1-p)^{k-1} (\frac{N}{N+1})^{k-1} F_{S_{k-1}}(W)}{P_{succ}(\Lambda^*)}
\end{equation}

\section{Numerical Results}
We validate our models using SimPy simulations. Table \ref{tab:results} summarizes representative data points.

\begin{table}[h]
\centering
\begin{tabular}{@{}llccr@{}}
\toprule
Case & Config & p & Theory & Sim \\ \midrule
1-Node & Mod & 0.2 & 1.45 & 1.43 \\
Tandem & N=4 & 0.1 & 2.91 & 2.91 \\
Feedback & N=10 & 0.1 & 31.61 & 31.80 \\ \bottomrule
\end{tabular}
\caption{Comparison of End-to-End Delay (Sojourn Time).}
\label{tab:results}
\end{table}

\begin{figure}[h]
    \centering
    \includegraphics[width=\columnwidth]{section_3_3_1_sojourn_vs_N_varying_p.png}
    \caption{Delay in Feedback Networks (Section 3.3.1).}
\end{figure}

\begin{figure}[h]
    \centering
    \includegraphics[width=\columnwidth]{section_3_3_2_tandem_delay_vs_N_varying_p.png}
    \caption{Delay in Tandem Networks (Section 3.3.2).}
\end{figure}

\section{Conclusion}
The proposed moment-matching and fixed-point frameworks provide a scalable solution for estimating performance in vulnerable networks.

\end{document}
